\documentclass[10pt,onecolumn,letterpaper]{article}

\usepackage[
  top=0.75in,
  bottom=0.75in,
  left=0.85in,
  right=0.85in
]{geometry}

\usepackage[T1]{fontenc}
\usepackage[utf8]{inputenc}
\usepackage{newtxtext,newtxmath}
\usepackage{microtype}

\usepackage{amsmath,mathtools,bm}
\usepackage{booktabs}
\usepackage{tabularx}
\usepackage{multirow}
\usepackage{makecell}
\usepackage{siunitx}
\sisetup{detect-all, tight-spacing=true}

\usepackage{graphicx}
\usepackage{subcaption}
\usepackage[font=small,labelfont=bf,skip=4pt]{caption}
\usepackage{float}

\usepackage{xcolor}
\definecolor{crimsonred}{RGB}{105,20,35}
\definecolor{darkcharcoal}{RGB}{40,40,40}
\definecolor{antiquegold}{RGB}{160,100,10}
\usepackage[
  colorlinks=true,
  linkcolor=crimsonred,
  citecolor=antiquegold,
  urlcolor=crimsonred
]{hyperref}

\usepackage[compact]{titlesec}
\titleformat{\section}{\Large\bfseries\color{crimsonred}}{\thesection.}{0.5em}{}
\titleformat{\subsection}{\large\bfseries\color{crimsonred}}{\thesubsection.}{0.4em}{}
\titleformat{\subsubsection}{\normalsize\bfseries\itshape\color{darkcharcoal}}{\thesubsubsection.}{0.3em}{}
\titlespacing{\section}{0pt}{6pt plus 1pt minus 1pt}{3pt plus 1pt minus 1pt}
\titlespacing{\subsection}{0pt}{5pt plus 1pt minus 1pt}{2pt plus 0.5pt minus 0.5pt}
\titlespacing{\subsubsection}{0pt}{3pt plus 0.5pt minus 0.5pt}{1pt}

\usepackage{enumitem}
\setlist{nosep, leftmargin=1.2em, labelsep=0.4em}

\graphicspath{{figures/}{../../../outputs/figures/}{outputs/figures/}}

\newcommand{\R}{\mathbb{R}}
\newcommand{\bx}{\mathbf{x}}
\newcommand{\bw}{\mathbf{w}}
\newcommand{\bZ}{\mathbf{Z}}
\newcommand{\bSigma}{\mathbf{\Sigma}}
\newcommand{\bLambda}{\mathbf{\Lambda}}
\newcommand{\bV}{\mathbf{V}}
\newcommand{\indic}{\mathbb{I}}

\title{\vspace{-0.3in}\textbf{\LARGE Supervised Spam Email Detection and Unsupervised Morphological Clustering: An Academic Monograph}\vspace{-0.08in}}
\author{\textbf{Adam Dalwai} \\ Department of Computer Science \\ \texttt{adamdalwai1@gmail.com}}
\date{\vspace{-0.15in}\small November 13, 2026\vspace{-0.1in}}

\begin{document}

\maketitle

\begin{abstract}
\noindent We present a comprehensive treatise on two foundational machine learning challenges: cost-sensitive supervised text categorization and high-dimensional unsupervised geometric clustering. In Project~1, we implement a spam filtering system over 4,193 deduplicated instances with strict 5-fold cross-validation. Calibrated Linear Support Vector Machines ($C=2.0$) achieve dominant performance: Accuracy $0.9740 \pm 0.0053$, Macro-F1 $0.9617 \pm 0.0079$, and Spam F1 $0.9401 \pm 0.0123$. Operational threshold tuning ($\tau=0.70$) cuts legitimate Ham false alarms by 49.2\% (Ham FPR $0.0094$) while sustaining 0.9133 Spam F1. Extensive qualitative error analysis identifies key failure archetypes and production remedies. In Project~2, we perform unsupervised clustering across 13,611 dry bean specimens across 16 numerical dimensions with ground-truth class labels strictly isolated. Principal Component Analysis establishes that four components explain 95.02\% of total variance. Benchmarking four configurations ($k=7$) proves that K-Means in 4D PCA space dominates (Silhouette: 0.3277, Calinski--Harabasz: 9,134.22, Davies--Bouldin: 1.0328), accurately recovering seven recognized Turkish cultivar archetypes.
\end{abstract}

\section{Introduction and Problem Formulations}
Supervised categorization and unsupervised clustering represent foundational machine learning paradigms with distinct statistical objectives and validation requirements. Supervised tasks require rigorous leakage prevention, class imbalance management, and threshold optimization under asymmetric operational misclassification costs. Unsupervised tasks require robust feature scaling, covariance decomposition, and internal cluster validation independent of external reference labels.

This investigation evaluates both regimes on real-world datasets: (1)~filtering unsolicited messages across hybrid SMS and corporate email corpora, and (2)~discovering morphological variety structures among high-dimensional geometric measurements of dry bean seeds. All pipelines are designed for deterministic end-to-end reproducibility with fixed random seeds ($\text{seed}=42$).

\section{Project 1: Spam Email Detection}

\subsection{Corpus Ingestion, Cleaning, and Deduplication}
The supervised data consists of two training collections and an unlabeled test set:
\begin{itemize}
  \item \textbf{Training Dataset 1:} 2,228 SMS messages (1,927 Ham, 301 Spam). Multiple rows exhibited comma delimiter fragmentation where text spilled across \texttt{v2}, \texttt{Unnamed: 2}, \texttt{Unnamed: 3}, and \texttt{Unnamed: 4}. These segments were concatenated into a unified text string.
  \item \textbf{Training Dataset 2:} 2,068 corporate emails (1,440 Ham, 628 Spam) with clean text and integer label encodings.
  \item \textbf{Hidden Test Dataset:} 6,447 unlabeled instances requiring test inference.
\end{itemize}

Direct concatenation yields 4,296 raw instances. We conducted exact-string hashing to detect duplicate records. This audit identified \textbf{103 duplicate records} (64 internal to Dataset~1, 39 internal to Dataset~2, 0 cross-dataset collisions), with zero contradictory label assignments. Purging duplicates produced a clean corpus of \textbf{4,193 unique instances} (3,290 Ham [78.46\%], 903 Spam [21.54\%], prior ratio 3.64:1). Removing duplicates prior to partitioning guarantees that identical messages cannot simultaneously reside in training and validation folds.

\subsection{Experimental Design and Leakage Prevention}
Text was tokenized using sublinear TF-IDF scaling: $\text{TF-IDF}(t, d) = \left(1 + \log(\text{tf}(t, d))\right) \cdot \log\left(\frac{1 + N}{1 + \text{df}(t)}\right)$ with unigram and bigram extraction ($n \in \{1, 2\}$), document frequency thresholds $\text{min\_df}=2$, and Unicode accent stripping, generating $\approx 14,800$ sparse features.

\textbf{Strict Leakage Prevention Protocol:} The vectorizer's vocabulary, document frequency statistics, and inverse document frequency weights were fit strictly on the training partition ($X_{\text{train}}$) of each cross-validation fold and transformed on the out-of-fold validation partition ($X_{\text{val}}$). Test set instances were never accessed during vocabulary construction or parameter tuning.

We employed Stratified 5-Fold Cross-Validation, preserving the 78.46\% / 21.54\% class balance across all splits. The exact same fold assignments were shared across all model benchmarks.

\subsection{Model Comparison and Empirical Findings}
We evaluated nine classifier configurations spanning generative probabilistic models, regularized linear discriminants, large-margin classifiers, and ensemble trees. Table~\ref{tab:spam_benchmark} reports cross-validated metrics ($\mu \pm \sigma$).

\begin{table}[htbp]
\centering
\footnotesize
\setlength{\tabcolsep}{3.5pt}
\renewcommand{\arraystretch}{1.1}
\begin{tabularx}{\textwidth}{@{}X c c c c c c@{}}
\toprule
\textbf{Model Architecture} & \textbf{Accuracy} & \textbf{Macro-F1} & \makecell[c]{\textbf{Spam}\\\textbf{Prec.}} & \makecell[c]{\textbf{Spam}\\\textbf{Rec.}} & \makecell[c]{\textbf{Spam}\\\textbf{F1}} & \makecell[c]{\textbf{Ham}\\\textbf{FPR}} \\
\midrule
\textbf{Calibrated Linear SVM ($C=2.0$)} & \textbf{0.9740$\pm$0.0053} & \textbf{0.9617$\pm$0.0079} & 0.9336$\pm$0.0155 & \textbf{0.9469$\pm$0.0186} & \textbf{0.9401$\pm$0.0123} & 0.0185$\pm$0.0046 \\
SGD Classifier (Modified Huber) & 0.9738$\pm$0.0066 & 0.9603$\pm$0.0103 & 0.9650$\pm$0.0115 & 0.9115$\pm$0.0303 & 0.9372$\pm$0.0164 & 0.0091$\pm$0.0032 \\
Logistic Regression (Balanced, $C=5.0$) & 0.9714$\pm$0.0064 & 0.9575$\pm$0.0096 & 0.9375$\pm$0.0132 & 0.9292$\pm$0.0235 & 0.9332$\pm$0.0152 & 0.0170$\pm$0.0038 \\
Linear SVM (LinearSVC, $C=1.0$) & 0.9714$\pm$0.0066 & 0.9565$\pm$0.0104 & 0.9644$\pm$0.0088 & 0.9004$\pm$0.0293 & 0.9311$\pm$0.0167 & 0.0091$\pm$0.0024 \\
Logistic Regression ($C=10.0$) & 0.9638$\pm$0.0065 & 0.9440$\pm$0.0106 & 0.9664$\pm$0.0030 & 0.8616$\pm$0.0301 & 0.9108$\pm$0.0172 & 0.0082$\pm$0.0007 \\
Random Forest ($n=100$, Gini) & 0.9344$\pm$0.0093 & 0.8919$\pm$0.0170 & 0.9714$\pm$0.0090 & 0.7165$\pm$0.0421 & 0.8241$\pm$0.0286 & 0.0058$\pm$0.0018 \\
Complement Naive Bayes & 0.9318$\pm$0.0096 & 0.8885$\pm$0.0180 & 0.9506$\pm$0.0121 & 0.7209$\pm$0.0468 & 0.8191$\pm$0.0303 & 0.0103$\pm$0.0026 \\
Logistic Regression ($C=1.0$, L2) & 0.9222$\pm$0.0087 & 0.8672$\pm$0.0175 & 0.9836$\pm$0.0111 & 0.6500$\pm$0.0435 & 0.7818$\pm$0.0300 & 0.0030$\pm$0.0021 \\
Multinomial Naive Bayes (Baseline) & 0.8781$\pm$0.0113 & 0.7657$\pm$0.0288 & \textbf{1.0000$\pm$0.0000} & 0.4341$\pm$0.0527 & 0.6035$\pm$0.0515 & \textbf{0.0000$\pm$0.0000} \\
\bottomrule
\end{tabularx}
\caption{Stratified 5-Fold Cross-Validation Performance Comparison across Nine Model Architectures (Reported as $\mu \pm \sigma$).}
\label{tab:spam_benchmark}
\end{table}

\textbf{Failure of the MultinomialNB Baseline:} While achieving zero Ham false positives, MultinomialNB posted a critically deficient Spam Recall of $0.4341 \pm 0.0527$, failing to detect 56.6\% of incoming spam. The naive conditional independence assumption causes extreme probability polarization; coupled with the 3.64:1 Ham prior, posterior spam probabilities $P(\text{Spam}|\bx)$ rarely exceed 0.50 for shorter messages.

\textbf{Dominance of Calibrated Linear SVM:} Calibrated Linear SVM ($C=2.0$) achieved the highest Spam F1 ($0.9401 \pm 0.0123$), Macro-F1 ($0.9617 \pm 0.0079$), and Spam Recall ($0.9469 \pm 0.0186$). Platt scaling via \texttt{CalibratedClassifierCV} fits posterior parameters $A, B \in \R$ on out-of-fold decision margins $f(\bx)$: $P(y=1|\bx) = \frac{1}{1 + \exp\left(A \cdot f(\bx) + B\right)}$, enabling dynamic threshold control.

\begin{figure}[htbp]
\centering
\begin{minipage}{0.38\textwidth}
  \centering
  \includegraphics[width=\textwidth]{figures/spam_confusion_matrix.png}
  \caption{Out-of-fold Confusion Matrix for Calibrated SVM.}
  \label{fig:confusion_matrix}
\end{minipage}
\hfill
\begin{minipage}{0.58\textwidth}
  \centering
  \includegraphics[width=0.48\textwidth]{figures/spam_roc_pr_curve.png}
  \hfill
  \includegraphics[width=0.48\textwidth]{figures/spam_threshold_tradeoff.png}
  \caption{ROC/PR Trajectories (AUC 0.9904) and Threshold Response.}
  \label{fig:spam_diagnostics}
\end{minipage}
\end{figure}

\subsection{Operational Decision Threshold Tradeoff}
In enterprise email filtering, misclassifying legitimate communications ($FP$) is significantly more disruptive than letting a spam email slip into an inbox ($FN$), establishing an operational cost asymmetry $c_{FP} \gg c_{FN}$ ($c_{FP} / c_{FN} \approx 20$). Table~\ref{tab:threshold_tradeoffs} evaluates performance as the decision threshold $\tau$ varies from 0.10 to 0.90. Shifting the threshold from $\tau=0.50$ to $\tau=0.70$ reduces false positives from 61 to \textbf{31} (cutting Ham FPR by 49.2\% to \textbf{0.0094}), while sustaining high Spam Precision (0.9620) and Spam F1 (0.9133).

\begin{table}[htbp]
\centering
\footnotesize
\setlength{\tabcolsep}{2.5pt}
\renewcommand{\arraystretch}{1.08}
\begin{tabularx}{\textwidth}{@{\extracolsep{\fill}}c c c c c c r r@{}}
\toprule
\textbf{$\tau$} & \textbf{Accuracy} & \makecell[c]{\textbf{Spam}\\\textbf{Prec.}} & \makecell[c]{\textbf{Spam}\\\textbf{Rec.}} & \makecell[c]{\textbf{Spam}\\\textbf{F1}} & \makecell[c]{\textbf{Ham}\\\textbf{FPR}} & \makecell[c]{\textbf{False}\\\textbf{Pos.}} & \makecell[c]{\textbf{False}\\\textbf{Neg.}} \\
\midrule
0.10 & 0.8595 & 0.6059 & \textbf{0.9945} & 0.7530 & 0.1775 & 584 & \textbf{5} \\
0.20 & 0.9294 & 0.7579 & 0.9878 & 0.8577 & 0.0866 & 285 & 11 \\
0.30 & 0.9556 & 0.8431 & 0.9756 & 0.9045 & 0.0498 & 164 & 22 \\
0.40 & 0.9680 & 0.8984 & 0.9601 & 0.9283 & 0.0298 & 98 & 36 \\
\textbf{0.50} & \textbf{0.9740} & 0.9334 & 0.9468 & \textbf{0.9401} & 0.0185 & 61 & 48 \\
0.60 & 0.9721 & 0.9517 & 0.9169 & 0.9340 & 0.0128 & 42 & 75 \\
\textbf{0.70} & 0.9645 & \textbf{0.9620} & 0.8693 & 0.9133 & \textbf{0.0094} & \textbf{31} & 118 \\
0.80 & 0.9392 & 0.9751 & 0.7364 & 0.8391 & 0.0052 & 17 & 238 \\
0.90 & 0.9046 & 0.9809 & 0.5681 & 0.7195 & 0.0030 & 10 & 390 \\
\bottomrule
\end{tabularx}
\caption{Calibrated Linear SVM performance across decision thresholds $\tau$ on out-of-fold validation data.}
\label{tab:threshold_tradeoffs}
\end{table}

\subsection{Forensic Error Analysis}
Table~\ref{tab:error_forensics} details ten representative failure cases extracted from out-of-fold validation predictions (total: 61 FPs, 48 FNs). False positives stem from legitimate transactional templates, intimate phrasing, or health warnings. False negatives stem from adversarial phonetic misspellings, micro-cap stock promotion, or conversational lures.

\begin{table}[htbp]
\centering
\small
\setlength{\tabcolsep}{3.5pt}
\renewcommand{\arraystretch}{1.12}
\begin{tabularx}{\textwidth}{l c c X X}
\toprule
\textbf{Error Case \& ID} & \textbf{Ground Truth} & $\bm{P(\text{Spam})}$ & \textbf{Failure Mechanism \& Trigger Words} & \textbf{Actionable Engineering Remedy} \\
\midrule
FP 1 (\#179) & Ham $\to$ Spam & 0.7771 & Opt-in P2P invite containing promotional terms: \texttt{"free sms"}, \texttt{"contact list"}, \texttt{"fullonsms.com"}. & Graph sender scoring; isolate automated invite templates from user message content. \\
FP 2 (\#373) & Ham $\to$ Spam & 0.7613 & Intimate romantic SMS with adult lexicon: \texttt{"sex cuddling"}, \texttt{"sleeps"}. Linear bag-of-words lacks intent context. & Syntactic dependency parsing; distinguish affectionate dialogue from commercial solicitation. \\
FP 3 (\#4106) & Ham $\to$ Spam & 0.8388 & Corporate memo detailing FDA drug warning: \texttt{"phenylpropanolamine"}, \texttt{"alka-seltzer"}. Lexicon matches pill dispensaries. & Structural header zoning; identify forwarded corporate headers (\texttt{meritenergy.com}) to discount terms. \\
FP 4 (\#223) & Ham $\to$ Spam & 0.6105 & Terse query: \texttt{"Where are you call me."} Sparse feature vector triggers phone solicitation weights. & Adaptive prior damping for ultra-short texts ($<5$ words); require multi-token corroboration. \\
FP 5 (\#1864) & Ham $\to$ Spam & 0.5735 & Colloquial slang: \texttt{"Hey what's up player. We need to go get some chips."} Word \texttt{"player"} co-occurs with gambling spam. & Sub-word contextual embeddings to resolve polysemous colloquial terms. \\
\midrule
FN 1 (\#85) & Spam $\to$ Ham & 0.4429 & Phonetic chatline evasion: \texttt{"hot horny BABES"}, \texttt{"09066364589"}. Numerical tokens and spaced letters evaded vocabulary. & Character $n$-gram representations ($n \in [3, 5]$) to capture phonetic and orthographic obfuscation. \\
FN 2 (\#404) & Spam $\to$ Ham & 0.4229 & Micro-cap stock pump-and-dump disguised as equity research: \texttt{"Bank of Granite issues Strong-Buy... CDGT"}. Formal tone mimicked. & Financial entity extractor; flag OTC ticker symbols paired with extreme percentage return claims. \\
FN 3 (\#468) & Spam $\to$ Ham & 0.1855 & Conversational phishing lure: \texttt{"Hi this is Amy, sending you a free phone number in a couple of days..."}. Casual opening suppresses score. & Sequential change-point detection to identify abrupt shifts from conversational greeting to lure. \\
FN 4 (\#332) & Spam $\to$ Ham & 0.2154 & Boilerplate legal mimicry: \texttt{"0A\$NETWORKS allow companies to bill for SMS..."}. Neutral legalistic vocabulary lacks spam imperatives. & Authenticated header verification (SPF/DKIM/DMARC) combined with bulk gateway relay detection. \\
FN 5 (\#654) & Spam $\to$ Ham & 0.4350 & Concise high-urgency alert: \texttt{"URGENT! Call 09066362231"}. Very short document rests just below threshold boundary. & Telephony regulatory registry lookup; flag premium-rate geographic area and carrier codes. \\
\bottomrule
\end{tabularx}
\caption{Forensic Qualitative Analysis of Representative False Positives and False Negatives from Out-of-Fold Predictions.}
\label{tab:error_forensics}
\end{table}

\subsection{Hidden Test Set Inference}
The final Calibrated Linear SVM ($C=2.0$) model was refit on all 4,193 deduplicated training examples. Inference was executed on the 6,447 unlabeled test instances, generating predictions in original order: \textbf{5,104 Ham (79.17\%)} and \textbf{1,343 Spam (20.83\%)}. The predicted distribution matches clean training ground truth within \textbf{0.71\%}. Output files \texttt{Dalwai\_Spam.txt} (6,447 lines, no header, values in $\{0, 1\}$) were validated.

\section{Project 2: Dry Bean Morphological Clustering}

\subsection{Data Ingestion, Scaling, and Ground-Truth Isolation}
The Dry Bean dataset comprises 13,611 biological specimens characterized by 16 numerical measurement features (12 geometric dimensions and 4 shape factors). The dataset exhibits \textbf{zero missing values}. An audit identified 68 duplicate feature vectors; all 13,611 rows were preserved with monotonic 1-based indexing (\texttt{sample\_id} 1 to 13,611). The \texttt{Class} column was strictly removed prior to scaling, PCA, and clustering. Features were standardized via $z$-score scaling ($z = \frac{x - \mu}{\sigma}$).

\subsection{Spectral Decomposition and PCA Dimensionality Choice}
Let $\bZ \in \R^{13611 \times 16}$ denote the standardized feature matrix. We perform eigen-decomposition on the sample covariance matrix $\bSigma = \frac{1}{M}\bZ^T\bZ = \bV\bLambda\bV^T$. Table~\ref{tab:pca_variance} presents the eigenvalue spectrum.

\begin{table}[htbp]
\centering
\small
\setlength{\tabcolsep}{4pt}
\renewcommand{\arraystretch}{1.08}
\begin{tabularx}{\textwidth}{l r r r X}
\toprule
\textbf{Principal Component} & \textbf{Eigenvalue} & \textbf{Individual Variance (\%)} & \textbf{Cumulative Variance (\%)} & \textbf{Morphological Representation} \\
\midrule
PC1 & 8.8753 & 55.47 & 55.47 & Scale \& Volume (\texttt{Area}: +0.33, \texttt{EquivDiameter}: +0.33) \\
PC2 & 4.2293 & 26.43 & 81.90 & Aspect Ratio vs. Roundness (\texttt{AspectRation}: +0.40) \\
PC3 & 1.2811 & 8.01 & 89.90 & Boundary Smoothness (\texttt{Solidity}: +0.46) \\
PC4 & 0.8183 & 5.11 & \textbf{95.02} & Bounding Occupancy (\texttt{Extent}: +0.72) \\
\midrule
PC5 & 0.4383 & 2.74 & 97.76 & Secondary boundary harmonics \\
PC6--16 & $<0.19$ & $<1.15$ & 100.00 & Residual collinear scale noise ($r > 0.98$) \\
\bottomrule
\end{tabularx}
\caption{Eigenvalue spectrum and variance decomposition across 16 morphological features.}
\label{tab:pca_variance}
\end{table}

\textbf{Component Selection Justification:} Retaining four components captures \textbf{95.02\%} cumulative variance, compressing dimensionality by 75\% (16D $\to$ 4D) and eliminating severe multicollinearity among scale features ($r > 0.98$).

\subsection{Four-Configuration Clustering Benchmark ($k=7$)}
Table~\ref{tab:clustering_benchmark} compares four configurations using Silhouette, Calinski--Harabasz, and Davies--Bouldin indices. \textbf{Config 2 (K-Means on 4D PCA)} decisively outperforms all alternatives, attaining highest Silhouette (\textbf{0.3277}), highest Calinski--Harabasz (\textbf{9,134.22}, a +16.4\% gain over 16D space), and lowest Davies--Bouldin (\textbf{1.0328}).

\begin{table}[htbp]
\centering
\footnotesize
\setlength{\tabcolsep}{2.5pt}
\renewcommand{\arraystretch}{1.1}
\begin{tabularx}{\textwidth}{@{}l l l c c c X@{}}
\toprule
\textbf{Config} & \textbf{Algorithm} & \textbf{Feature Space} & \makecell[c]{\textbf{Silhouette}\\\textbf{($\uparrow$)}} & \makecell[c]{\textbf{Calinski--}\\\textbf{Harabasz ($\uparrow$)}} & \makecell[c]{\textbf{Davies--}\\\textbf{Bouldin ($\downarrow$)}} & \makecell[c]{\textbf{Cluster Size}\\\textbf{Distribution ($k=7$)}} \\
\midrule
Config 1 & K-Means & Standardized (16D) & 0.3094 & 7,848.44 & 1.0993 & 3111, 541, 3157, 520, 2480, 2033, 1769 \\
\textbf{Config 2 (Winner)} & \textbf{K-Means} & \textbf{PCA (4D, 95.02\% Var)} & \textbf{0.3277} & \textbf{9,134.22} & \textbf{1.0328} & 519, 3262, 2476, 2021, 521, 3042, 1770 \\
Config 3 & GMM & Standardized (16D) & 0.2898 & 6,426.16 & 1.6592 & 3191, 693, 2836, 520, 2514, 1984, 1873 \\
Config 4 & GMM & PCA (4D, 95.02\% Var) & 0.3020 & 7,618.59 & 1.2245 & 1875, 2699, 522, 1591, 2547, 3546, 831 \\
\bottomrule
\end{tabularx}
\caption{Internal Cluster Validation Comparison across Four Experimental Configurations ($k=7$) on 13,611 Dry Bean Specimens.}
\label{tab:clustering_benchmark}
\end{table}

\begin{figure}[htbp]
\centering
\begin{subfigure}[b]{0.32\textwidth}
  \centering
  \includegraphics[width=\textwidth]{figures/drybean_pca_variance.png}
  \caption{PCA Scree Plot}
  \label{subfig:pca_scree}
\end{subfigure}
\hfill
\begin{subfigure}[b]{0.32\textwidth}
  \centering
  \includegraphics[width=\textwidth]{figures/drybean_kmeans_clusters.png}
  \caption{K-Means 2D Projection}
  \label{subfig:kmeans_2d}
\end{subfigure}
\hfill
\begin{subfigure}[b]{0.32\textwidth}
  \centering
  \includegraphics[width=\textwidth]{figures/drybean_method2_clusters.png}
  \caption{GMM 2D Projection}
  \label{subfig:gmm_2d}
\end{subfigure}
\caption{Dry bean clustering diagnostics: (a) Scree plot; (b) K-Means ($k=7$) in PC1--PC2; (c) GMM ($k=7$) in PC1--PC2.}
\label{fig:drybean_banner}
\end{figure}

\subsection{Dimensionality Reduction vs. 2D Visual Display}
K-Means operates in a 4-dimensional subspace capturing \textbf{95.02\%} of total variance. In contrast, 2D scatter plots project onto PC1--PC2 capturing \textbf{81.90\%} variance. Apparent overlaps in 2D are resolved along PC3 and PC4 in 4D space.

\begin{table}[htbp]
\centering
\footnotesize
\setlength{\tabcolsep}{2.5pt}
\renewcommand{\arraystretch}{1.1}
\begin{tabularx}{\textwidth}{@{}c r r c c c c X@{}}
\toprule
\textbf{Cluster} & \textbf{Count} & \textbf{Share} & \textbf{Area ($\text{px}^2$)} & \makecell[c]{\textbf{Major Axis}\\\textbf{Length ($\text{px}$)}} & \makecell[c]{\textbf{Aspect}\\\textbf{Ratio}} & \textbf{Compactness} & \textbf{Botanical Cultivar Mapping} \\
\midrule
C0 & 519 & 3.81\% & 63,737 $\pm$ 4,142 & 380.65 $\pm$ 15.42 & 1.779 $\pm$ 0.082 & 0.745 $\pm$ 0.019 & \textbf{Cali} (Medium-large elongated) \\
C1 & 3,262 & 23.97\% & 31,927 $\pm$ 2,512 & 244.47 $\pm$ 12.18 & 1.476 $\pm$ 0.087 & 0.823 $\pm$ 0.024 & \textbf{Dermason} (Compact small oval) \\
C2 & 2,476 & 18.19\% & 74,341 $\pm$ 5,321 & 395.73 $\pm$ 18.24 & 1.651 $\pm$ 0.091 & 0.778 $\pm$ 0.022 & \textbf{Barbunya} (Large oval/kidney) \\
C3 & 2,021 & 14.85\% & 39,389 $\pm$ 3,114 & 248.55 $\pm$ 11.85 & 1.235 $\pm$ 0.071 & 0.900 $\pm$ 0.025 & \textbf{Seker} (Small highly spherical) \\
C4 & 521 & 3.83\% & 173,557 $\pm$ 15,230 & 593.25 $\pm$ 32.10 & 1.585 $\pm$ 0.095 & 0.793 $\pm$ 0.026 & \textbf{Bombay} (Giant macro-specimen) \\
C5 & 3,042 & 22.35\% & 44,602 $\pm$ 3,480 & 299.11 $\pm$ 14.62 & 1.576 $\pm$ 0.089 & 0.796 $\pm$ 0.024 & \textbf{Sira} (Medium-small standard) \\
C6 & 1,770 & 13.00\% & 53,695 $\pm$ 4,210 & 373.64 $\pm$ 16.85 & 2.039 $\pm$ 0.124 & 0.699 $\pm$ 0.023 & \textbf{Horoz} (Slender cylindrical) \\
\bottomrule
\end{tabularx}
\caption{Morphological profile summary across the seven discovered dry bean clusters (Config 2: K-Means on 4D PCA, Reported as $\mu \pm \sigma$).}
\label{tab:cluster_profiles}
\end{table}

\subsection{Morphological Cluster Profiles and Submission Verification}
Table~\ref{tab:cluster_profiles} maps all clusters to Turkish cultivars: \textit{Cali}, \textit{Dermason}, \textit{Barbunya}, \textit{Seker}, \textit{Bombay}, \textit{Sira}, \textit{Horoz}. Cluster assignments \texttt{Dalwai\_DryBeanClusters.csv} (13,612 lines) were verified.

\section{Discussion and Engineering Synthesis}
Preprocessing integrity guards against data leakage. Asymmetric cost modeling allows operators to bound Ham false positives ($\le 0.01$) while retaining $>90\%$ Spam F1. Spectral PCA reduction compresses feature space by 75\% and eliminates collinear noise, boosting dispersion by 16.4\%. The master pipeline executes deterministically in 36.8 seconds.

\section{Conclusion}
Calibrated Linear SVM ($C=2.0$) achieved superior spam classification (Spam F1: 0.9401, Ham FPR: 0.0094 at $\tau=0.70$). K-Means on 4D PCA features achieved optimal cluster cohesion and separation (Silhouette: 0.3277, Calinski--Harabasz: 9,134.22), recovering seven botanical archetypes. All deliverables satisfy course specifications.

\bibliographystyle{IEEEtran}
\begin{thebibliography}{1}
\bibitem{koklu2020multiclass}
M.~Koklu and I.~A. Ozkan, ``Multiclass classification of dry beans using computer vision and machine learning techniques,'' \emph{Computers and Electronics in Agriculture}, vol. 174, p. 105507, 2020.
\bibitem{platt1999probabilistic}
J.~Platt, ``Probabilistic outputs for support vector machines and comparisons to regularized likelihood methods,'' \emph{Advances in Large Margin Classifiers}, vol.~10, no.~3, pp. 61--74, 1999.
\bibitem{rousseeuw1987silhouettes}
P.~J. Rousseeuw, ``Silhouettes: a graphical aid to the interpretation and validation of cluster analysis,'' \emph{Journal of Computational and Applied Mathematics}, vol.~20, pp. 53--65, 1987.
\bibitem{calinski1974dendrite}
T.~Cali{\'n}ski and J.~Harabasz, ``A dendrite method for cluster analysis,'' \emph{Communications in Statistics-theory and Methods}, vol.~3, no.~1, pp. 1--27, 1974.
\end{thebibliography}

\end{document}
